% ============================================================================================
% RQ3 — Through what process do skill inconsistencies cost the agent?
% 441 paired cells (49 tasks x 3 models x trials r1-r3, RQ2 cell-selection rule).
% ============================================================================================

\section{RQ3: Through what process do skill inconsistencies cost the agent?}
\label{sec:rq3}

RQ2 measures \emph{how much} a repaired skill changes outcomes; it cannot say \emph{why}. A repair
changes one element of one skill, and the aggregate effect is the sum of many executions in which an
agent may or may not encounter that element, may or may not adopt it, and may or may not recover. This
section opens that aggregate along a three-stage chain --- \emph{exposure}, \emph{response},
\emph{outcome} --- asking, in order,

\begin{itemize}
  \item[\textbf{RQ3.1}] \emph{Exposure.} Do agents come into contact with the skill at all?
  \item[\textbf{RQ3.2}] \emph{Response.} Given contact, what do they do with the inconsistent element?
  \item[\textbf{RQ3.3}] \emph{Outcome.} What does each response cost, in score and in compute?
\end{itemize}

The questions are ordered by dependence: a response presupposes contact, and a cost presupposes a
response.

\paragraph{Trajectory coding.} We analyse the trajectories of the 441 paired cells used by RQ2 (49
tasks $\times$ 3 models $\times$ trials r1--r3), each contrasting \textsc{OriginalSkill} against
\textsc{FixedSkill} on an identical task. Cells are resolved with RQ2's selection rule --- current
task-package image, latest study among duplicates --- because 134 of the 882 (axis, task, arm, trial)
keys have between two and five candidate runs spread across studies; resolving by manifest order
silently mixes stale versions and, as we report, inverts one of our principal results. For each task
the repair diff yields three task-specific anchors: the files the repair changed, code-like tokens
present only in the shipped skill (\emph{shipped-only}), and tokens present only in the repaired skill
(\emph{fixed-only}); the diff also yields the verifier checks tied to the changed behaviour. A
deterministic coder then walks each trajectory, separating what the agent \emph{authored} (tool
arguments) from what tools \emph{returned} (file contents, command output), and records skill contact,
use of a changed file, manifestation of a shipped-only anchor, per-check verifier failures, and
turn-level timestamps. Two decisions carry the analysis. First, \emph{contact is not use}: reading a
file that contains an anchor is not adopting it, and running the skill's script is not triggering the
inconsistency; only an authored shipped-only form, or execution of a changed file, counts as use.
Second, \emph{response is attributed by pairing}: a behaviour is attributed to the inconsistency only
when the \textsc{FixedSkill} arm of the same cell does not exhibit it, which separates the
inconsistency's effect from a model's habits. The coder is rule-based; its behavioural labels are
validated on a stratified sample by hand, reported in \S\ref{sec:threats}.

\begin{table}[t]
\centering\small
\caption{Trajectory coding scheme. Exposure is contact; the response is what follows, attributed
against the paired repaired arm. Outcome is measured by the harness and verifier, never inferred.}
\label{tab:coding}
\begin{tabular}{@{}p{2.6cm}p{12.1cm}@{}}
\toprule
Code & Operational definition \\
\midrule
\textbf{Exposed} & the agent reads the skill, whether by opening a file under \texttt{skills/} or by
  loading it through the CLI's \texttt{Skill} tool \\
\textbf{Unused} & the agent never reads the skill; the inconsistency is outside the execution \\
\textbf{Bypass} & exposed, but the inconsistent element is never used: no changed file is executed and
  no shipped-only form is authored \\
\textbf{Follow} & the inconsistent behaviour is used and reaches the artifact: a tied verifier check
  fails in \textsc{OriginalSkill} and passes in \textsc{FixedSkill} \\
\textbf{Recover} & the inconsistent behaviour fires, but the run still passes the tied checks \\
\textbf{BothFail} & the tied check fails in both arms; the repair was not decisive (Section~3.3) \\
\bottomrule
\end{tabular}
\end{table}

% --------------------------------------------------------------------------------------------
\subsection{RQ3.1: Exposure --- contact is nearly universal}
\label{sec:rq3-exposure}

\begin{table}[t]
\centering\small
\caption{Exposure. A run is exposed if the agent reads the skill at all; the remaining 13 runs are
unused (8 where neither arm read the skill, 5 where only \textsc{OriginalSkill} skipped it). The route
columns overlap: \emph{path} is opening a file under \texttt{skills/}, \emph{tool} is loading the skill
by name through the CLI's \texttt{Skill} tool.}
\label{tab:exposure}
\begin{tabular}{@{}lrrrrrr@{}}
\toprule
 & & & & \multicolumn{3}{c}{Access route} \\
\cmidrule(l){5-7}
Model & $n$ & Exposed & Rate & Path only & Tool only & Both \\
\midrule
Claude Haiku 4.5  & 147 & 144 &  98\%  &  75 & 21 & 48 \\
Qwen3.6-35B       & 147 & 137 &  93\%  & 120 & 1  & 16 \\
DeepSeek-V4-Flash & 147 & 147 & 100\%  & 119 & 1  & 27 \\
\midrule
All               & 441 & 428 & 97.1\% & 314 & 23 & 91 \\
\bottomrule
\end{tabular}
\end{table}

Contact is the norm, not the exception (Table~\ref{tab:exposure}): across all 441 executions the agent
reads the skill in 97.1\% of runs, and the strongest model reads it in every single run. Only 13 runs
(2.9\%) are unexposed. Two methodological points guard this measure. Contact is counted on both access
routes, because Haiku reaches the skill through the \texttt{Skill} tool far more often than the other
two (69 runs against 28 and 17); counting only path-based access would understate exposure by five
points overall and by fifteen points for Haiku. And the measure is not inflated by directory listing:
no execution inspects \texttt{skills/} without also reading a file.

\textbf{Finding 1.} Whether an agent sees the inconsistency is essentially constant across models and
therefore cannot explain the RQ2 gaps. Contact is not where the models differ.

% --------------------------------------------------------------------------------------------
\subsection{RQ3.2: Response --- what differs is what happens after contact}
\label{sec:rq3-response}

\begin{table}[t]
\centering\small
\caption{Response, given exposure. \emph{BothFail} is where the tied check fails under both skills
and is analysed further in Section~3.3.}
\label{tab:response}
\begin{tabular}{@{}lrrrr@{}}
\toprule
Response & Haiku & Qwen & DeepSeek & All \\
\midrule
Recover $\rightarrow$ success   & 40 & 33 & \textbf{82} & 155 \\
Follow $\rightarrow$ failure   & \textbf{44} & 35 & 18 &  97 \\
BothFail                       & 40 & 43 & 20 & 103 \\
Used, no effect                &  9 & 11 & 19 &  39 \\
Bypass                         & 11 & 15 &  8 &  34 \\
Unused                         &  3 & 10 &  0 &  13 \\
\bottomrule
\end{tabular}
\end{table}

The models are indistinguishable in exposure but sharply distinguishable in what they do next
(Table~\ref{tab:response}). Haiku follows the inconsistency into failure in 30\% of its runs, Qwen in
24\%, and DeepSeek in 12\%; recovery shows the reverse ordering, at 27\%, 22\% and 56\%. Haiku and
Qwen also bypass the inconsistent element more often than DeepSeek, and Qwen is the only model that
leaves skills unread at a non-trivial rate.

\textbf{Finding 2.} The models differ after contact, not at contact: whether an agent adopts what it
reads, rather than whether it reads at all, is the decisive step.

\textbf{Finding 3.} The per-model counts of following-into-failure --- Haiku 44, Qwen 35, DeepSeek 18
--- reproduce the ordering of the RQ2 repair gains exactly (Haiku $+0.136$, Qwen $+0.085$, DeepSeek
$+0.023$, not significant). The aggregate repair effect is therefore the number of executions in which
the agent adopts the inconsistency and delivers it. This ordering is visible only under RQ2's
cell-selection rule: resolving duplicate cells by manifest order yields Haiku 34 and Qwen 35,
spuriously inverting it.

\begin{figure*}[t]
\centering
\includegraphics[width=\textwidth]{fig-response-type.pdf}
\caption{Response composition by inconsistency type and model ($n = 36, 21, 15, 15, 12$ executions
per type respectively; bars are normalised to 100\%, so segments are read as shares). \textsc{BothFail}
is a task-difficulty control. The models diverge after contact in every type, but most sharply in
\texttt{over\_promised}, where the share followed into failure runs from 22\% (DeepSeek) to 50\%
(Haiku); Qwen converts the same type into the unrecoverable control most often (25\%). Models are
ordered as in Table~\ref{tab:rq2-effectiveness}.}
\label{fig:response-type}
\end{figure*}

Figure~\ref{fig:response-type} disaggregates the same responses by model and by inconsistency type.
The model gradient is not the product of a single type: it holds within every type, and it is
steepest in \texttt{over\_promised}, where the share followed into failure ranges from 22\% (DeepSeek)
to 50\% (Haiku).

\textbf{Finding 4.} Following the inconsistency into failure is concentrated in \texttt{over\_promised}
(41 of the 97 follow cases), in which the documentation promises a capability the code never
implements. Conversely, \texttt{semantic\_conflict} is recovered in 13 of its 18 cases: a contradiction
in meaning is noticed and worked around far more often than a missing capability.

% --------------------------------------------------------------------------------------------
\subsection{RQ3.3: Outcome --- the score is lost in one place, the compute everywhere}
\label{sec:rq3-outcome}

\begin{table}[t]
\centering\small
\caption{Outcome by response. $n$ is the number of paired executions and \emph{fires} the subset in
which the inconsistency actually manifested, on which the footprint is measured. $\Delta$score and net
$\Delta$token are \textsc{FixedSkill} $-$ \textsc{OriginalSkill} per run, so a positive token value
means the shipped skill spent more. The footprint is the segment from the turn the incident first
fired to the turn the agent settled on the corrected form, or to the end of the run if it never did.
\textsc{BothFail} is analysed further in Table~\ref{tab:bothfail}; the \textsc{Bypass} and
\textsc{Unused} rows are not attributable because the inconsistent element was never used.}
\label{tab:outcome}
\begin{tabular}{@{}lrrrrr@{}}
\toprule
Response & $n$ (fires) & $\Delta$score & net $\Delta$token & footprint turns & footprint (s) \\
\midrule
Follow $\rightarrow$ failure  & 97 (87)  & $+0.366$ & $+34$k & 14.4 & 104 \\
Recover $\rightarrow$ success & 155 (155) & $-0.038$ & $-57$k & 12.3 & 60 \\
Bypass                        & 34 (4)   & $+0.097$ & $+2$k  &  3.8 & 40 \\
BothFail                      & 103 (87) & $+0.030$ & $+7$k  & 19.6 & 169 \\
Used, no effect               & 39 (3)   & $-0.047$ & $-44$k & 11.3 & 24 \\
Unused                        & 13 (3)   & $+0.127$ & $+4$k  & 16.0 & 88 \\
\bottomrule
\end{tabular}
\end{table}

\textbf{Finding 5.} The repair's entire benefit is concentrated in one cell of the matrix
(Table~\ref{tab:outcome}). Following the inconsistency into failure costs $0.366$ of the score and
repairing it recovers it; every other class sits at $\Delta \approx 0$. Moreover the per-incident loss
is nearly independent of the model: among the follow executions, $\Delta$score is $+0.351$ (Haiku,
$n{=}44$), $+0.372$ (Qwen, $n{=}35$) and $+0.393$ (DeepSeek, $n{=}18$). Models differ in how often
an inconsistency is followed, not in what it costs when it is.

\textbf{Finding 6.} Footprint and net cost diverge, and the table shows this within single rows.
Recovering and following produce similar footprints (12.3 and 14.4 turns), yet the net token
difference is $-57$k for the former and $+34$k for the latter: on Recover cells the repaired arm
spends $57$k tokens \emph{more}, because a more complete skill induces more work, which offsets the
detour the agent would otherwise take. Repair buys correctness, not thrift. Footprint is also not
cross-model comparable in seconds, because the local 4-bit Qwen deployment is an order of magnitude
slower (a follow detour costs it 205 seconds against 49 for Haiku and 27 for DeepSeek at similar turn
counts); turn counts, not seconds, are the comparable footprint.

\begin{table}[t]
\centering\small
\caption{The \textsc{BothFail} class is not a single control: it decomposes by how many tied checks
the repair changed. A quarter of it is genuinely neutral, but more than a third shows a real gain and
more than a quarter is actively harmful.}
\label{tab:bothfail}
\begin{tabular}{@{}lrrr@{}}
\toprule
Repair's effect on the tied checks & $n$ & share & mean $\Delta$score \\
\midrule
Failure count reduced (partial fix)  & 38 & 37\% & $+0.352$ \\
Failure count unchanged              & 37 & 36\% & $+0.011$ \\
Failure count increased (adverse)    & 28 & 27\% & $-0.380$ \\
\midrule
Total                                & 103 & 100\% & \\
\bottomrule
\end{tabular}
\end{table}

\textbf{Finding 7.} The executions in which the tied check fails under both skills are not a task
difficulty floor; they are heterogeneous (Table~\ref{tab:bothfail}). In 38 the repair reduced the
number of failed tied checks and recovered $+0.352$ of score --- a genuine, partial success. In 37 the
repair changed nothing, which is the true floor. And in 28 executions --- 6\% of all paired runs ---
repairing the skill made the outcome worse by $-0.380$ on average. These adverse cases cluster in a
handful of tasks (e.g.\ \texttt{promql-lint-gate}, \texttt{pdf-forms}, \texttt{dingtalk-todo}) and
reproduce the negative-$\Delta$ tasks already visible in RQ2; they show that a repair, while removing
one inconsistency, can still mislead the agent.

\textbf{Finding 8.} Attribution is only possible where the treatment was administered. In
\textsc{Bypass} and \textsc{Unused} the original arm never used the inconsistent element, so their
score deltas are not the inconsistency's cost; the \emph{fires} column makes this visible ($4$ of $34$
and $3$ of $13$). \textsc{Bypass} measures avoidance, not harm.

% --------------------------------------------------------------------------------------------
\subsection{RQ3 summary}
\label{sec:rq3-summary}

\begin{table}[t]
\centering\small
\caption{RQ3 in one table: the aggregate repair gain of RQ2 decomposes into the number of executions
in which the agent adopts the inconsistency and delivers it.}
\label{tab:rq3-summary}
\begin{tabular}{@{}lrrr@{}}
\toprule
 & Haiku & Qwen & DeepSeek \\
\midrule
Exposure rate                    & 98\% & 93\% & 100\% \\
Follow $\rightarrow$ failure     & \textbf{44} & 35 & 18 \\
Recover $\rightarrow$ success    & 40 & 33 & \textbf{82} \\
RQ2 repair gain $\Delta$(F$-$O)  & $+0.136$ & $+0.085$ & $+0.023$ \\
\bottomrule
\end{tabular}
\end{table}

Exposure is near-universal and nearly constant across models, so it cannot explain the RQ2 gaps. What
differs is the response: the strongest model recovers from the inconsistency twice as often as it
follows it, while the weaker two follow it about as often as they recover, and the count of
follow-into-failure executions orders the models exactly as RQ2's repair gains do
(Table~\ref{tab:rq3-summary}). The cost of an inconsistency has two separable parts --- the score lost
when it is followed, and the compute consumed by the detours of runs that recover --- and the second
leaves no trace in the aggregate score.
